\section{操作自救：一整套 ADHD 原生方法库}

认知打完了地基，这一章是整份讲义的操作核心。这里的每一个方法，都不是「通用自律技巧」的换皮，而是\term{顺着 ADHD 算法本身设计出来的}——它们不要求你提高意志力，而是要求你降低意志的需求。

先给一个大原则：\term{顺着算法，不逆着算法}。所有方法都可以归结为三句话——

\begin{checklist}
  \item \textbf{以广促深}：广度算法便宜，那就让它便宜着干活，别逼它走深度；
  \item \textbf{以被动代主动}：主动切换贵，被动切换便宜，那就多制造被动，少依赖主动；
  \item \textbf{以环境代意志}：意志耗能贵，环境布局便宜，那就把意志活干的事交给环境。
\end{checklist}

\subsection{方法一：恰当场 / 抽象轨道}

这是整套方法的底座（BV11ucZzxEkx）。核心动作只有两个：

\term{第一，把一堆事情，在内心合并成「一件事」。}你这学期有十本书要读，那这十本书就是一件事：「我这学期该读的东西」。A 项目、B 汇报、C 考试，凡是你觉得「同时符合当下和长远」的，都并进一个\term{恰当场}，然后告诉自己：\term{我今天只做一件事}。

\term{第二，内部随便跳，但不跳出轨道。}念头从 A 跳到 B 再跳回 C，无所谓——外部看你在频繁切换，内部看你一直在这条宽轨道上运行。这避开了最贵的东西：主动任务切换（那意味着抑制当前冲动）。你跳的是被动的、便宜的、自发的念头。

范围怎么定？\term{一开始要定得大一些}，大到「下至小学数学、上至前沿研究」都装得下的程度。不用担心因此整天做简单事——ADHD 大脑本身追求新奇，简单的东西无聊得更快，它自己会往深处走。

\begin{insight}
一句话：\term{把「必须做 ABCD」的排序，换成「我只在这一片场里随便蹦」}。排序是高耗能的抑制过程，蹦是低耗能的自发过程。让大脑觉得「我是自由的」，但自由的范围恰好是你定的。
\end{insight}

\subsection{方法二：程式性计划：计划做 A，不限定内容}

ADHD 完全可以做计划，但要做\term{程式性计划}（BV16VckzqEMp）：计划「下午做 A」，但不限定任何具体内容。

于是下午就变成：随便翻 A 的资料、书、视频，翻到哪个有兴致就停一下。\term{只要翻得足够多，激活兴趣的概率总是存在的}——哪怕一下午没写出什么「正事」，你肯定翻到了某个让你眼睛一亮、顺手写了两笔的点。

这里有个心理坎必须过：你怕这样太散、不系统，怕「正常人一下午按部就班看了很多，而我一直在乱看，差距越来越大」。原视频的回答非常直接：\term{不用担心，我在清华待了这么多年，知道正常人一下午也学不了多少}。线性积累模型的错误，我们第六章已经拆过了。

注意一个要点：\term{A 的范围要足够大}，大到你可以随便跳。计划「下午学编程」而不是「下午学第一章」，这样跳到数学、跳到工具、跳到相关想法都不算出轨。

\subsection{方法三：毛点计划：计划只是锚点}

ADHD 做计划，最常见的失败方式是：把计划当执行清单，做不完就自责。原视频给了一个更聪明的定位（BV1Jj8vzWEAQ）：\term{计划只是头脑里的毛点（锚点），不是执行路径}。

你有一个计划，就像心里种了一个钩子。之后在日常里，你会不断布置环境、制造刺激，让一些「啪的一下」的想法冒出来——\term{一拍脑袋冒出来的想法，往往正好和你的计划有关}。这时候立刻去行动。

为什么「啪的一下」很重要？因为\term{这个想法冒出来的时候，你的多巴胺是够的}。启动困难的大脑，最缺的就是这一下。原视频甚至回忆自己早年太严肃地对待这类想法——总想着「这想法靠不靠谱、别人做过没有」，结果把最有动力的启动机会全否定了，退回枯燥的读文献轨道，越想越不想动，陷入恶性循环。

\begin{checklist}
  \item 计划 = 毛点：想得足够远，能激活就算成功，不要求绝对可控；
  \item 环境 = 触发器：文件分类、桌面布局、收藏夹，都朝着「让相关想法冒出来」布置；
  \item 行动 = 追着「啪的一下」走：冒出来就做，别自我否定。好的想法，常常不是在按部就班里迭代出来的，是在兴奋的尝试里突然冒出来的。
\end{checklist}

\subsection{方法四：框架欺骗：撤掉「必须」警报}

这是对付「排异反应」的经典招数（BV1WBPezFEqE）。

你可能有这种体验：心里一旦确定「接下来两小时我要做 A」，身体立刻产生抵抗——想睡、想刷手机、想逃跑。原视频的解释是：\term{「必须做 A」这句话，在神经层面是一个警报}——系统预判「即将进入高耗能模式」，于是提前启动防御，表现为抗拒、拖延、犯困。

解法是\term{框架欺骗}：告诉自己「接下来两小时干啥都可以」，但实际行为局限在 A 相关范围内——你随便翻 A 的资料，只是心里没有「必须」。警报撤了，防御关了，大脑进入开放状态；一旦某个点恰好点燃了兴趣，系统就会\term{带着兴趣自动倒向 A}。

\begin{insight}
这个策略的精妙之处：\term{你给了大脑自由的幻觉，却没有偏离目标}。自由状态下的大脑更容易被激活，而激活的那个点，正好在 A 的范围里。用自由骗过防御，用范围保住目标。
\end{insight}

\subsection{方法五：布局神经元：环境布局代替意志力}

这是原视频的另一个核心概念（BV1MWwJzxEh5，BV13HoGYiEgF）：\term{主动注意力你控制不了，但你可以布置环境，改变你的「猎物偏好」}。

具体操作：

\begin{checklist}
  \item 最重要的一篇文章 / 文件，单独占一个屏幕，占满，别让一堆窗口跟你抢视线；
  \item 正在写的记录一直开着，随时可以落笔；
  \item 关键资料放在「走神后三秒内能切到」的位置，而不是藏在文件夹深处；
  \item 反过来，容易带偏你的东西，从视野里拿走——降低它们的暴露频率，就等于降低它们在你注意力竞争里的权重。
\end{checklist}

原理：你没法命令自己「想到正事」，但可以让「正事」在你的视野里反复出现，于是念头在竞争时更可能落到它上面。这就是\term{以环境代意志}——你改不了注意力，但你能改环境给注意力喂什么。

\subsection{方法六：冲动引导成网，随手记录}

第三章讲过冲动理性。这里给它的操作化（BV1apQ2ByEN3）：

\term{冲动来了，别硬压，也别全放。}硬压，抑制耗能爆表；全放，方向完全失控。正确姿势是：\term{记录它，然后继续}。冒出一个想法，用 Obsidian 随手记两行（甚至只是几个自然语言的关键词），清掉一部分工作记忆，然后跳回原来的事。这个想法是种子，之后它可能自己长成深刻的问题。

关键细节：\term{首次冲动大多没啥用，但你别懒得记}。因为二次冲动、三次冲动往往是从首次冲动的残留上长出来的，越记越长、越记越深。原视频原话：\term{「记忆不得想法，就是这些首次冲动上长出来的二次冲动」}——很多好问题，就是靠「随手记一下」喂出来的。

方向怎么控？不是靠压，是靠\term{调节环境暴露概率}：你同时想提高 A 的比例、降低 B 的比例，那就减少环境里引发 B 想法的刺激，多布置引发 A 想法的刺激。冲动已经来了，就释放掉；没冲动的时候，就做 A，并布置更利于 A 自然激发的环境。

\subsection{方法七：隐形回血：管理分心，而非消除分心}

第三章讲过隐形回血。这里补操作细节（BV16hxrzYEpT）：

分心分两类，区别对待：

\begin{checklist}
  \item \textbf{无新内容的冲动}（突然想刷手机、开无关网页）：用身份认同压下去——「我是一个气度沉稳、能按顺序行事的人」，这一句就有用；
  \item \textbf{有价值的创新性想法}（冒出一个新的连接、一个猜想）：立刻记下来或问 AI——这是回血，是好想法的种子，也是启动门槛的降低剂。
\end{checklist}

为什么这算回血？因为全神贯注的效率，会被事后累抵消——硬撑两小时敲完代码，之后躺三小时；允许有价值的分心，四个小时敲完，中间产出一堆想法，还不太累。\term{后者是更划算的}。别当电动车的司机，却告诉车「你要跑得更快，所以不准充电」。

\subsection{方法八：唤醒度三档管理 + 假性疲劳处理}

第二章讲过唤醒度三档。这里把它变成一个可执行的日常流程（BV16mXnBxEE1，BV1MWwJzxEh5）：

\begin{checklist}
  \item \textbf{高唤醒}：做线性任务、写重要的东西，趁热打铁；
  \item \textbf{中唤醒}：随手翻找——电脑里、书里、Obsidian 里，看看有什么能激发兴趣；或做点本身不需要高唤醒的事；
  \item \textbf{低唤醒}：躺下，允许大脑自己放想法。用手机里「方便思考且你觉得有意义」的内容，或者直接问 AI 几个问题，重新唤醒。别硬撑，也别内疚——这不算偷懒，这是这套系统的重启方式。
\end{checklist}

处理假性疲劳的口诀是：\term{想法上高自由度，操作上低自由度}。

\begin{checklist}
  \item \textbf{想法上高自由度}：告诉自己「今天不一定要写这个，随便偏出去写点别的，也是成果」——撤掉「必须」警报；
  \item \textbf{操作上低自由度}：要写的文件放桌面最显眼处，记录一直开着——让任何想法都能在几秒内转向它。
\end{checklist}

\subsection{方法九：音乐策略}

「ADHD 学习能不能听音乐」——原视频给了比「能/不能」更细的答案（BV1T1XSBnE6w）：

\begin{checklist}
  \item \textbf{机械熟练任务}（整理、搬运、抄录、打码）：放音乐。不放音乐时用来防分心的自我监控成本，可能比放音乐还高；音乐（尤其是纯音乐）还能占住想跳出的念头，收窄跳跃范围；
  \item \textbf{随意翻找内容}：放不放影响不大；
  \item \textbf{看似机械、但需要预测每一步的任务}（不熟的步骤）：关音乐。因为你的大脑每一步都在预测，那需要主动注意力，别和音乐抢带宽。
\end{checklist}

核心逻辑：机械任务上争分夺秒没意义——它们永远只是准备工作，关键在「做」，不在「快」。把省下来的意志留给真正重要的任务。

\subsection{方法十：模糊分类：给广度搜索装一个范围}

ADHD 的搜索是广度优先的（BV1xpLFzaECh）。当头脑里缺乏「粗糙的分类」，搜寻范围过大，大脑就会「一团浆糊」——看理工科内容尤其容易「脑内温暝、像烧开水一样」。解法是\term{模糊分类}：

\begin{checklist}
  \item 给一本书的十三四章，自己概括成四五个大类，按你的方式起名字（「攻防法服」那种级别的粗糙分类）；
  \item 再往上，给整个学科建大目录——根据你读的文章和脑里的想法，把学科分块；
  \item 分类的名字必须概括特征，不能瞎起；但边界模糊没关系，模糊正是它的优点。
\end{checklist}

为什么模糊比严格好？\term{严格化要你平白多记很多东西，还破坏了强弱顺序}。而粗糙分类能一眼看出「这两个东西同属一类，所以能用同样的方法处理」——省下的算力是巨大的。古人把东西分成阴阳五行、地水火风，对古人来说可能真的有用，因为它限定了搜寻范围。

\subsection{方法十一：慕容水战法（或战或亡）}

这是原视频最戏剧性的一招（BV1zDpkzZEtX）。先说结论：\term{慢性压力是 ADHD 的天敌，急性压力反而释放认知资源}。

慢性压力（做不完要挨骂、要失业、做不出来）推高皮质醇，皮质醇一高，启动锁死——你陷入「这任务做完也没啥用 + 做不完要被评判」的死循环，五点就想睡。而\term{无路可退的急性压力}，是另一套机制：当大脑只剩下「或战或亡」一个选项，所有「做不做、什么时候做、怎么做更好」的内耗瞬间消失，目标变得无比清晰——就是眼前这件事。

原视频把它叫\term{慕容水战法}：想象自己像慕容垂那样，七十岁、久病缠身，但这一次出征没有退路——「此刻只负责把眼前这一仗打完，不再想后续」。这个框架一旦进入，认知资源反而被释放，内耗消失，反而能干活了。

\begin{notehot}
必须强调：\term{这招不能常用}。它有点「先伤己、再伤人」的性质，是伤敌一千自损八百的绝地求生。它的价值在于：当你被慢性压力困死的时候，知道还有这么一把钥匙——把「想太多」换成「只打眼前这一仗」，哪怕只是假想地进入这个状态，也能达到部分效果。自我欺骗只要能起实际的好效果，从神经学角度看就是真实。
\end{notehot}

\subsection{方法十二：工具层}

原视频专门推荐过一套工具组合（BV1xMonY7EGD，BV16hxrzYEpT）：

\begin{checklist}
  \item \textbf{Obsidian}：打开文件快、可折叠、有关系图谱——天生适合「想法冒出来随手记、跳回来再找」的工作流；
  \item \textbf{虚拟桌面划分}：桌面一放最重要的一篇文章（单内容占满），桌面二放 Obsidian，桌面三放常用查询网站。主动切过去是自发的、可以的；被 Alt+Tab 弹出一堆窗口被动干扰，才是要避免的；
  \item \textbf{AI 即时反馈}：低唤醒 / 清醒睡眠状态时，你的符号语言处理能力掉线，但自然语言还在。把数学对象、猜想、想法用自然语言讲给 AI，当下就拿到反馈——原视频认为这对 ADHD 极其友好，是过去没有的「即时奖励」。
\end{checklist}

\subsection{方法十三：游戏与刺激的重新管理}

最后是一点生活管理（BV1U1XbBfEtE）。ADHD 经常在游戏上栽跟头，但原视频的建议很反直觉：

\term{不要在一天内规划「游戏两小时」。}这种精确配额带来的结果往往是：玩得不够尽兴、一超标就觉得「反正超标了，干脆玩到吐」，变成一天十几个小时、连续几个月。改成\term{直接滑一两个整天去玩}——把游戏当成一件被允许的事，而不是要偷摸超标的事。意志耗能降下来，反而容易收手。

对睡眠的态度也一样（BV16mXnBxEE1）：不必纠结「必须马上睡着」。晚上到点躺下就行，睡不着时头脑也在休息；躺床上用手机问 AI、看有趣内容，不叫没休息——对 ADHD 来说，那就是休息的一种。

\subsection{一个通用的启动流程}

把以上十二个方法串起来，就是一条可复制的启动流水线：

\begin{enumerate}[leftmargin=2em,itemsep=2pt]
  \item 先处理无聊：机械任务放音乐，中低唤醒就翻找激发，别硬坐；
  \item 用框架欺骗撤掉「必须」警报：告诉自己「接下来两小时干啥都行」，实际只翻 A 相关；
  \item 在恰当场里被动切换：念头冒出来就跟，跳出轨道的不跟；
  \item 有价值的想法：随手记录或问 AI，那是回血和种子；
  \item 假性疲劳：动态调自由度——想法上放开，操作上把文件放最显眼处；
  \item 被慢性压力困死时：换一次慕容水战法，只打眼前这一仗。
\end{enumerate}

\key{这一章没有一个方法是「逼自己」的。它们全是同一句话的变体：别硬改大脑，去改它运行的场地。}
